\section{Conclusions}\label{sec:conclusion}

\subsection{Answers to the four questions}

\textbf{What distinguishes the three call types?} They differ along the frequency axis, not in
duration or bandwidth. Peak frequency separates the groups after correction ($H=12.28$,
$p=0.013$), and the median first fundamental decreases from flat (298 Hz) through modulated
(287 Hz) to two-voice (274.5 Hz). The two-voice call additionally carries a second fundamental,
which is a structural feature rather than a continuous measurement. A classifier using only the
shared continuous features reaches a macro-$F_1$ of 0.458; admitting the structural indicators
raises it to 0.952, but that gain reflects the definition of the call types, not a better
acoustic description.

\textbf{How far does the signal travel?} Fidelity is nearly constant to 32 m and collapses
between 64 and 256 m. At the operational threshold $\tau=0.5$ the effective ranges are 84.7 m
(flat), 144.6 m (modulated) and 169.7 m (two-voice). The ordering is robust across thresholds and
across most weightings, but not under an envelope-heavy weighting --- at $\tau=0.7$ with weight
0.5 on the envelope correlation the two-voice radius collapses to 12.1 m --- nor across model
families: at $\tau=0.5$ a four-parameter Weibull reverses the modulated and two-voice order. The
three correlations are not independent evidence; the spectrogram and amplitude-spectrum metrics
are strongly correlated (per-type $r=0.81$--$0.97$).

\textbf{Do penguins respond, and to what?} Yes, strongly, and the paired design is what makes
this visible: 64.8\% against 12.9\%, McNemar OR 37.0 (95\% CI 9.9--311). Response declines with
distance (OR 0.49 per $e$-fold) and with group size (OR 0.94 per bird), and the three call types
differ from control (OR 13.7--22.4) but not clearly from one another. The distribution of the
eight response combinations changes under playback ($p=0.0002$); whether call type shifts the
probability of a specific combination is not established ($p=0.137$, 58\% of expected counts
below five). Time to first response has a conditional median of about 2 s and lengthens by a
factor of 1.27 per $e$-fold of distance.

\textbf{What is the communication-effective range?} With
$\mathrm{CEI}(d)=F(d)\times P_{\mathrm{recruit}}(d)$, the behaviour-limited half-effectiveness
radii are 137 m (flat), 175 m (modulated) and 203 m (two-voice) on the empirical curve, with the
two-voice value falling to 158 m under a logistic fit. Recruitment probability itself declines
only weakly with distance (OR 0.87 per $e$-fold, $p=0.59$), so the radius is driven mainly by
fidelity rather than by a behavioural gradient. If a 3 dB increase were equivalent to the assumed
0.10 drop in fidelity, the radii would fall to 102, 141 and 171 m. We translate these into a
precautionary 200 m buffer with a 3 dB limit in the 0.5--2 kHz band.

\subsection{Future work}

Three extensions would materially improve the conclusions. First, recording the observation
window explicitly would turn $W=60$ s from an assumption into a measurement. Second, marking
individual birds would make $\rho$ identifiable and would separate group-level from
individual-level response. Third, an experiment that links a playback trial to a propagation
measurement of the same call would make the perception term $q(d)$ estimable rather than
imported, and would close the largest gap in the present analysis.
